%%%%%%%%%%%%%%%%%%%%%%%%%%%%%%%%%%%%%%%%%%%%%%%%%%%%%%%%%%%%%%%%%%%%%%%%%%%
% Función: capítulo de ejemplo de metodología de la tesis por compendio.
%
% Uso: copia de restauración; normalmente no debes modificarla.
% Tras restaurarla, modifica la copia activa para adaptarla a tu trabajo.
%%%%%%%%%%%%%%%%%%%%%%%%%%%%%%%%%%%%%%%%%%%%%%%%%%%%%%%%%%%%%%%%%%%%%%%%%%%

\ifthenelse{\equal{\myLanguage}{english}}
{
  \chapter{Methodology}
}
{
  \chapter{Metodología}
}
\label{cha:compendium-methodology}

\ifthenelse{\equal{\myLanguage}{english}}
{
  \section{Research Design}
  \subsection{General approach}
}
{
  \section{Diseño de la investigación}
  \subsection{Enfoque general}
}

\ifthenelse{\equal{\myLanguage}{english}}
{
  Describe the common methodology connecting the publications, including the experimental design, data, analytical procedures and validation criteria. Refer to the included publications for details while keeping this chapter understandable as a coherent account of the thesis methodology.
}
{
  Describe la metodología común que conecta las publicaciones, incluidos el diseño experimental, los datos, los procedimientos de análisis y los criterios de validación. Remite a las publicaciones incluidas para los detalles, manteniendo este capítulo como una explicación coherente de la metodología de la tesis.
}

\ifthenelse{\equal{\myLanguage}{english}}
{
  \subsection{Data and experimental protocol}

  Explain how the data are obtained, prepared and divided, and identify the variables and controls used throughout the publications.

  \section{Evaluation Methodology}
}
{
  \subsection{Datos y protocolo experimental}

  Explica cómo se obtienen, preparan y dividen los datos, e identifica las variables y los controles utilizados en el conjunto de las publicaciones.

  \section{Metodología de evaluación}
}

\ifthenelse{\equal{\myLanguage}{english}}
{
  Figure~\ref{fig:compendium-example} and Table~\ref{tab:compendium-example} are simple examples that you can replace with material from your thesis.
}
{
  La figura~\ref{fig:compendium-example} y la tabla~\ref{tab:compendium-example} son ejemplos sencillos que puedes sustituir por material de tu tesis.
}

\begin{figure}[htbp]
  \centering
  \includegraphics[width=0.55\textwidth]{Figure1}
  \ifthenelse{\equal{\myLanguage}{english}}
  {
    \caption{Example figure in the extended summary.}
  }
  {
    \caption{Ejemplo de figura en el resumen extendido.}
  }
  \label{fig:compendium-example}
\end{figure}

\begin{table}[htbp]
  \centering
  \renewcommand{\arraystretch}{1.2}
  \ifthenelse{\equal{\myLanguage}{english}}
  {
    \caption{Example comparison of experimental results.}
  }
  {
    \caption{Ejemplo de comparación de resultados experimentales.}
  }
  \label{tab:compendium-example}
  \begin{tabular}{|l|c|c|}
    \hline
    \ifthenelse{\equal{\myLanguage}{english}}{Method}{Método} &
    \ifthenelse{\equal{\myLanguage}{english}}{Metric A}{Métrica A} &
    \ifthenelse{\equal{\myLanguage}{english}}{Metric B}{Métrica B} \\
    \hline
    \ifthenelse{\equal{\myLanguage}{english}}{Baseline}{Referencia} & 0.72 & 0.68 \\
    \hline
    \ifthenelse{\equal{\myLanguage}{english}}{Proposed method}{Método propuesto} & 0.84 & 0.79 \\
    \hline
  \end{tabular}
\end{table}

%%% Local Variables:
%%% TeX-master: "../../book"
%%% End:
